\chapter{System Design and Methodology}

\section{Architecture}

\begin{figure}[h!]
\centering
\fbox{\includegraphics[width=15cm,height=9cm]{healthcare_fraud_workflow.png}}
\caption{Proposed System Architecture}
\label{fig:healthcare_architecture}
\end{figure}

The proposed system, ``AI-Based Healthcare Claims Fraud Detection with Medical Validation and Explainable AI'', is designed as a data-driven decision support framework for identifying clinically inconsistent patterns and potentially fraudulent healthcare providers. The system consists of two complementary analytical components: medical validation using NFHS-5 health survey data and machine learning-based potential fraud detection using Medicare healthcare claims data.

The first component uses NFHS-5 data to identify medically inconsistent combinations of reported health conditions, treatment status, medication usage, and glucose measurements. Rule-based validation is applied to identify suspicious clinical patterns. These patterns are treated as indicators requiring further investigation and are not considered confirmed cases of fraud.

The second component uses Medicare healthcare claims data to perform provider-level potential fraud detection. The claims, beneficiary, inpatient, and outpatient datasets are processed to generate provider-level features. These features include total claims, total reimbursement, average reimbursement, number of unique patients, inpatient claim statistics, and outpatient claim statistics.

The extracted provider-level features are used to train supervised machine learning models including Random Forest, Logistic Regression, and XGBoost. The models learn patterns associated with providers labelled as potentially fraudulent or non-fraudulent in the Medicare dataset.

The predictions generated by the machine learning models are evaluated using classification metrics such as accuracy, precision, recall, F1-score, ROC-AUC, and Average Precision. In addition, different classification thresholds are analyzed to study the trade-off between identifying more potentially fraudulent providers and reducing false positive predictions.

Finally, Explainable AI is incorporated using SHAP (SHapley Additive exPlanations). SHAP provides information about the contribution of individual features to model predictions, thereby improving the interpretability and transparency of the fraud detection process.

\section{Dataset and Data Processing}

The proposed framework uses two independent datasets for complementary purposes.

\subsection{NFHS-5 Dataset}

The National Family Health Survey-5 (NFHS-5) dataset contains health-related information collected across India during the 2019--2021 survey period. The dataset contains approximately 724,115 female respondents and a large number of health and demographic variables.

For the proposed system, clinically relevant variables related to diabetes, hypertension, treatment, medication, and glucose measurements are selected. The selected variables are processed in chunks to efficiently handle the large dataset.

The preprocessing stage includes handling invalid or missing survey codes and selecting valid observations for the medical validation rules. The resulting data is then passed through the rule-based validation system.

\subsection{Medicare Claims Dataset}

The Medicare dataset consists of provider, beneficiary, inpatient claim, and outpatient claim information. The provider-level ``PotentialFraud'' label is used as the target variable for supervised machine learning.

The claims data is aggregated at the provider level to create a compact representation of provider activity. The resulting dataset contains 5,410 providers, including providers labelled as potentially fraudulent and non-fraudulent.

\section{Medical Validation using NFHS-5}

The medical validation component identifies combinations of health information that may indicate clinical inconsistency.

The following validation rules are applied:

\begin{itemize}

\item Diabetes reported without corresponding treatment.

\item Hypertension reported without corresponding treatment.

\item Treatment reported without a corresponding diabetes diagnosis.

\item Treatment reported without a corresponding hypertension diagnosis.

\item High glucose measurement without a reported diabetes diagnosis.

\item Reported diabetes without prescribed medication.

\end{itemize}

Each rule generates a binary indicator representing whether the corresponding inconsistency is present.

A weighted risk score is then calculated using the following formulation:

\[
RiskScore =
3R_1 + 3R_2 + 5R_3 + 5R_4 + 4R_5 + 2R_6
\]

where each $R_i$ represents one of the identified medical validation rules.

The resulting scores are divided into different risk categories:

\begin{itemize}
\item Clean: Risk score = 0
\item Low: Risk score $\leq$ 2
\item Medium: Risk score $\leq$ 4
\item High: Risk score $\leq$ 6
\item Critical: Risk score $>$ 6
\end{itemize}

These categories provide a mechanism for prioritizing records that contain potentially inconsistent medical information.

\section{Feature Engineering for Medicare Claims}

The Medicare claims data is transformed into provider-level features before applying machine learning models.

The main features generated include:

\begin{itemize}

\item \textbf{Total Claims}: Total number of claims associated with a provider.

\item \textbf{Total Reimbursement}: Total reimbursement amount associated with the provider.

\item \textbf{Average Reimbursement}: Average reimbursement per claim.

\item \textbf{Unique Patients}: Number of distinct beneficiaries associated with the provider.

\item \textbf{Inpatient Claims}: Number of inpatient claims associated with the provider.

\item \textbf{Inpatient Reimbursement}: Total reimbursement associated with inpatient claims.

\item \textbf{Inpatient Average Reimbursement}: Average reimbursement for inpatient claims.

\item \textbf{Outpatient Claims}: Number of outpatient claims associated with the provider.

\item \textbf{Outpatient Reimbursement}: Total reimbursement associated with outpatient claims.

\item \textbf{Outpatient Average Reimbursement}: Average reimbursement for outpatient claims.

\end{itemize}

These features provide a statistical representation of provider-level healthcare activity and are used as inputs to the supervised learning models.

\section{Machine Learning Models}

Three supervised machine learning algorithms are used for provider-level potential fraud detection.

\subsection{Random Forest}

Random Forest is an ensemble learning algorithm that combines multiple decision trees to produce a classification result. In the proposed system, a Random Forest classifier with 100 trees and balanced class weights is used.

The use of class weighting helps address the imbalance between potentially fraudulent and non-fraudulent providers in the dataset.

\subsection{Logistic Regression}

Logistic Regression is used as a baseline linear classification model. StandardScaler is applied to normalize the numerical features before training the model.

Class weighting is also incorporated to reduce the effect of class imbalance.

\subsection{XGBoost}

XGBoost is a gradient boosting-based machine learning algorithm used for classification. It builds an ensemble of decision trees sequentially, with each subsequent tree attempting to improve the errors made by previous trees.

The model uses parameters such as 100 estimators, maximum tree depth of 4, learning rate of 0.1, subsampling, and column subsampling. Class imbalance is handled using the calculated positive-class weight.

\section{Model Training and Evaluation}

The provider-level Medicare dataset is divided into training and testing sets using an 80:20 stratified split. Stratification ensures that the class distribution is maintained across the training and testing sets.

Five-fold Stratified Cross-Validation is also performed to obtain a more reliable estimate of model performance.

The models are evaluated using the following metrics:

\begin{itemize}

\item Accuracy

\item Precision

\item Recall

\item F1-score

\item ROC-AUC

\item Average Precision

\end{itemize}

Since fraud detection is generally an imbalanced classification problem, precision and recall are particularly important for understanding model behaviour.

\section{Classification Threshold Analysis}

In addition to evaluating the default classification threshold, different probability thresholds are examined for the Random Forest model.

Thresholds such as 0.5, 0.4, 0.3, and 0.2 are evaluated to understand how changing the decision boundary affects precision, recall, and the number of detected potentially fraudulent providers.

Lowering the threshold generally allows the system to identify more positive cases, but may also increase the number of false positive predictions. Therefore, threshold analysis provides an additional mechanism for selecting an operating point according to the requirements of the investigation process.

\section{Explainable AI using SHAP}

Explainable Artificial Intelligence is incorporated into the proposed system using SHAP.

SHAP assigns contribution values to individual features for a model prediction. These values help explain which provider-level characteristics contributed towards a particular prediction.

For the Random Forest model, SHAP TreeExplainer is used to calculate feature contribution values. The resulting explanations can be visualized using SHAP summary and feature importance plots.

The use of SHAP improves the transparency of the machine learning system by allowing investigators to understand the factors associated with a model's prediction rather than relying solely on a binary classification output.

\section{Decision Support Framework}

The final system is intended to function as a decision-support mechanism rather than an automated fraud adjudication system.

The overall workflow can be summarized as:

\begin{enumerate}

\item Collect healthcare and health survey data.

\item Preprocess and validate the input data.

\item Apply medical consistency rules to NFHS-5 records.

\item Aggregate Medicare claims at the provider level.

\item Generate provider-level statistical features.

\item Train supervised machine learning models.

\item Evaluate model performance using multiple classification metrics.

\item Analyze different classification thresholds.

\item Generate SHAP-based explanations for model predictions.

\item Prioritize suspicious records or potentially fraudulent providers for further human investigation.

\end{enumerate}

The combination of rule-based medical validation, supervised machine learning, and Explainable AI provides a structured approach for healthcare fraud analysis while keeping final investigation and decision-making under human review.\begin{figure}
    \centering
    \includegraphics[width=0.5\linewidth]{healthcare_fraud_workflow.png}
    \caption{Enter Caption}
    \label{fig:placeholder}
\end{figure}
